\section{Finite actions and
extensions}\label{finite-actions-and-extensions}

\textbf{Theorem UD.Orbits (orbit data and normalizers).} Let a finite
group \(A\) act on a finite set \(\Omega\), and choose one
representative \(H\) of each conjugacy class of subgroups that occurs as
a point stabilizer. Let \(m_H\) be the number of orbits isomorphic to
\(A/H\), with \(m_H=0\) for all omitted classes. Then the numbers
\(m_H\) determine the \(A\)-set up to equivariant isomorphism, and \[
\ker(A\curvearrowright\Omega)=\bigcap_{m_H>0}\operatorname{core}_A(H).
\tag{5.3}
\] Its equivariant permutation group is abstractly \[
C_{\operatorname{Sym}(\Omega)}(A)\cong
\prod_H\bigl((N_A(H)/H)\wr\operatorname{Sym}(m_H)\bigr).
\tag{5.4}
\] If the action \(\rho:A\hookrightarrow\operatorname{Sym}(\Omega)\) is
faithful, conjugation gives an exact sequence \[
1\longrightarrow C_{\operatorname{Sym}(\Omega)}(\rho(A))
\longrightarrow N_{\operatorname{Sym}(\Omega)}(\rho(A))
\longrightarrow
\{\alpha\in\operatorname{Aut}(A):m_{\alpha(H)}=m_H\text{ for all }H\}
\longrightarrow1.
\tag{5.5}
\] No splitting of this sequence is asserted.

\emph{Proof.} Every orbit is equivariantly \(A/H\), where \(H\) is the
stabilizer of a selected point, and two such transitive sets are
isomorphic exactly when the stabilizers are conjugate. This proves the
classification by multiplicities and the kernel formula.

An equivariant permutation preserves the isomorphism type of each orbit.
For one orbit \(A/H\), an equivariant map is determined by the image of
\(H\). That image \(nH\) must be fixed by \(H\), and bijectivity makes
its stabilizer equal to \(H\); hence \(n\in N_A(H)\). Conversely each
such \(n\) gives a permutation \(aH\mapsto anH\). Composition
corresponds to multiplication in the opposite order, so inversion
identifies this permutation group with \(N_A(H)/H\). For \(m_H\) copies,
one may independently choose such a map in every copy and permute the
copies, yielding the wreath product. Different isomorphism types are
independent, proving (5.4).

Faithfulness makes conjugation by a normalizing permutation \(\tau\)
induce a unique automorphism \(\alpha\) of \(A\), with \[
\tau\rho(a)\tau^{-1}=\rho(\alpha(a)).
\] If a point has stabilizer \(H\), its image under \(\tau\) has
stabilizer \(\alpha(H)\); therefore the multiplicities must be
invariant. The kernel of this conjugation map is the centralizer.
Conversely, suppose \(\alpha\) preserves the multiplicities. Pair each
orbit with a target orbit whose stabilizer class is \(\alpha(H)\), and
choose base points \(x,y\) with actual stabilizers \(H,\alpha(H)\),
adjusting within the target orbit if needed. Define
\(\tau(ax)=\alpha(a)y\). The stabilizer equality makes this well-defined
and bijective on that pair of orbits, and the paired definitions combine
to a permutation of \(\Omega\). It satisfies the displayed intertwining
equation, proving surjectivity and exactness. \(\square\)

\textbf{Proposition UD.ExtensionOrbits (classification above a retained
base).} Let a group \(B\) act on sets \(Z,E\), and let \(p:Z\to E\) be
equivariant. For one representative \(e\) of each orbit in
\(B\backslash E\), there is a natural bijection \[
B\backslash Z\ \cong\ \coprod_{e\in B\backslash E}B_e\backslash p^{-1}(e).
\tag{5.6}
\]

\emph{Proof.} Translate any \(z\in Z\) so that \(p(z)\) is the selected
representative of its orbit. If two translated points over that
representative are \(B\)-equivalent, an element taking one to the other
fixes their image \(e\), so belongs to \(B_e\). Conversely a
\(B_e\)-equivalence is already a \(B\)-equivalence. These two
observations prove surjectivity and injectivity of (5.6). \(\square\)
